\chapter{Configuration \& Scaling}

\section{Scaling Laws}
LLM performance scales predictably with:
\begin{itemize}
    \item \textbf{N:} Number of parameters
    \item \textbf{D:} Amount of training data (tokens)
    \item \textbf{C:} Compute budget
\end{itemize}

A common heuristic (Chinchilla scaling) suggests that for every parameter, you should train on approximately 20 tokens. For a 1B parameter model, you need $\approx$ 20B tokens.

\section{Model Presets}
The `pure_transformer` provides several configurations optimized for A100 training.

\begin{table}[h]
\centering
\begin{tabular}{|l|c|c|c|c|}
\hline
\textbf{Config} & \textbf{Params} & \textbf{Layers} & \textbf{Hidden} & \textbf{Heads} \\ \hline
Tiny & 45M & 12 & 512 & 8 \\ \hline
Small & 125M & 16 & 768 & 12 \\ \hline
Medium & 350M & 24 & 1024 & 16 \\ \hline
Medium-Large & 400M & 20 & 1152 & 16 \\ \hline
Large & 760M & 24 & 1536 & 16 \\ \hline
Extra Large & 1.5B & 32 & 1920 & 20 \\ \hline
\end{tabular}
\caption{Model Configurations in pure\_transformer}
\end{table}

\begin{lstlisting}[language=Python, caption=Configuration Dataclass]
# configs/model_config.py

@dataclass
class TransformerConfig:
    vocab_size: int = 50304
    hidden_size: int = 1024
    intermediate_size: int = 2816  # approx 2.75x hidden for SwiGLU
    num_layers: int = 24
    num_heads: int = 16
    num_kv_heads: int = 4  # GQA ratio 4:1
    head_dim: int = 64
    max_seq_length: int = 2048
\end{lstlisting}

\section{Conclusion}
This handbook has covered the essential components of a modern decoder-only transformer. By understanding these foundations—Tokenization, RoPE, GQA, SwiGLU, and robust Training Loops—you are well-equipped to explore, modify, and train your own LLMs using the `pure_transformer` codebase.
